\honest{Diseased Thinking: Dissolving Questions about Disease}{Scott Alexander}{2010}

I open with a quotation from George Will: treating bad behavior as caused rather than chosen strips a person ``of personhood, and moral dignity.''

Sandy is obese, and fictional. Her husband tells her to stop eating like a pig. Her doctor says obesity is mostly genetic and suggests a pill and surgery. Her sister says it is a valid way to live, and that prejudice against fat people is like racism. They all quarrel: the husband says the others excuse her and sap her willpower, the doctor says the husband ignores the real causes and the sister legitimizes a danger, and the sister calls the husband a jerk and the doctor a servant of Big Pharma. Similar quarrels surround alcoholism, depression, attention deficit disorder, social anxiety disorder (``didn't we used to call this `shyness'?''), oppositional defiant disorder, homosexuality and more. I focus on the husband and the doctor; whether a condition is a disease decides, for most people, which of them is right.

``Disease'' is like the made-up word ``blegg'': a label for a set of traits that usually come together. In a factory that sorts red, smooth, vanadium-filled cubes from blue, furry, palladium-filled eggs, ``blegg'' for the eggs is useful because the traits are correlated; if blue things were as likely to be vanadium cubes, the word would tell you nothing. I list six: a cause of the kind biology studies, such as genes or viruses; being involuntary, ``completely immune to the operations of free will''; being rare; being unpleasant; a clear line between two populations, not a normal distribution; and treatment with chemicals or radiation. Cancer has all six. So do heart attacks, the flu and diabetes. \nb{The common kind of diabetes is diagnosed by a cutoff on a continuous blood test, with a middle band called prediabetes, so it is not sharply divided from health.} Dwarfism fails the fifth, aging the third, and homosexuality the fourth, since it is not necessarily unpleasant. Obesity arguably meets the first, fourth and sixth, but hardly the second, third and fifth. So asking whether obesity is ``really'' a disease is like asking whether Pluto is really a planet.

Once we say which criteria it meets, there is no further fact. But the word matters because of hidden inferences: it decides whether the person gets sympathy, and whether they may seek medical treatment. Cancer patients are always ``brave''; the shy are told to man up, while people with social anxiety disorder get sympathy. Treatment for a non-disease is a ``quick fix,'' and some doctors I have talked to hesitate to suggest gastric bypass even when indicated, thinking it wrong to use medicine for a character issue. If we can decide those two directly, we will have ``unasked'' the question.

On sympathy, I describe the usual view. Bad choices come from free will, ``a spiritual entity,'' so blame depends on whether a condition is biological or spiritual. A peaceful man whose brain tumor drove him to a killing spree, and who is peaceful again once it is removed, is not blamed. But no one knows how to tell whether a condition like depression is spiritual or biological. My reply: ``Determinist consequentialists can do better. We believe it's biology all the way down.'' Brain tumors and poor taste in music are both biological, but taste is open to social influences such as praise, condemnation and introspection, and tumors are not. On this view no one deserves bad treatment for its own sake (``Saddam Hussein doesn't deserve so much as a stubbed toe''), but punishment can have good consequences: hurting bank robbers prevents robberies, and condemning alcoholism may make it less attractive. My rule: ``if giving condemnation instead of sympathy decreases the incidence of the disease enough to be worth the hurt feelings, condemn; otherwise, sympathize.'' I admit that the rule rests ``on philosophy that the majority of the human race would disavow,'' but say it gives intuitively right results. Yelling at a cancer patient for letting cells divide will not cure the cancer; telling a lazy person to get up and work might cure the laziness.

\nb{The test in the rule is whether a condition responds to social pressure. That is close to the second criterion for disease, being ``completely immune to the operations of free will,'' restated without free will. The post had promised to decide sympathy ``independently of the central node `disease' or of the criteria that feed into it.''}

For laziness, condemnation ``very well might'' work; for obesity, the case I began with, the answer may depend on the person, on who does the condemning, and on what other treatments exist. \nb{Whether condemnation reduces a condition is a question of fact. The rule speaks of the ``incidence'' of a condition, and the post endorses condemnation that deters others, so the rule allows condemning a person who cannot change, as long as others are deterred. The post does not discuss that case.}

On treatment: may people use a drug when a social remedy also works? My answer: ``Of course, why wouldn't it be?'' Under the free-will view, a biological fix for a spiritual problem may seem dehumanizing, or a band-aid; to someone for whom it is biology all the way down, it is not. Others say easy fixes stop people learning responsibility, which shows a bias toward the way things are, as a reversal test reveals. If personal responsibility mattered more than not being addicted, they would want schoolchildren addicted to heroin so they could learn to quit. Then I quote Katja Grace: people who made sacrifices for a cause oppose easier solutions, because the sacrifice has become a sign of virtue. People who diet and exercise hard and then campaign against a pill that makes it easy show ``some of the same principles.'' \nb{Grace's pattern was about sacrifices for problems ``that don't threaten them personally.'' Obesity does threaten them personally.} I grant that drugs raise reasonable objections, such as side effects, cost, dependence, overhyped efficacy and patients who cannot consent, and say that these apply to cancer drugs too and have the same good-enough solutions. People who want effective treatment should not be denied it or stigmatized for seeking it.

``We should blame and stigmatize people for conditions where blame and stigma are the most useful methods for curing or preventing the condition.'' \nb{The post has not shown, for any condition, that they are.}
